\section{As-Built Microstructure and Mechanical Performance}
\label{sec:micro}

\subsection{Characteristic AM microstructures of HEAs}

Across powder-bed processes, as-built HEAs display fine cellular or dendritic solidification substructures, high dislocation densities arranged in cell walls, and epitaxially grown columnar grains aligned with the build direction; texture, cell size, and phase balance are all tunable through laser power, scan speed, hatch spacing, and scan strategy \cite{cagirici2024,han2025}. In single-phase FCC alloys such as CoCrFeMnNi, LPBF produces a homogeneous solid solution with grain sizes well below cast equivalents, which raises yield strength via the Hall--Petch effect while retaining ductility \cite{lpbf_cantor_2024}. In eutectic AlCoCrFeNi$_{2.1}$, SLM produces a fine, relatively homogeneous FCC/B2 duplex structure whose inter-phase composition spread is far smaller than in cast or annealed material \cite{slm_ehea_2025}. WAAM deposits are coarser and more anisotropic: elongated columnar grains, dendritic substructure, and reheated interpass regions \cite{andersson2026waam,liu2021waam}.

\subsection{Defects and their engineering meaning}

The defect population of AM HEAs mirrors that of AM metals generally: lack-of-fusion pores, gas porosity, solidification cracks, and inclusions, plus HEA-specific risks such as hot cracking in Al-rich compositions and elemental loss by evaporation \cite{cagirici2024}. For marine service the important point is that every pore or crack intersecting a wetted surface is a candidate pit or crevice initiation site, and internal porosity can connect into crevice-like networks after surface machining \cite{nass2025marine}. Hot isostatic pressing closes internal porosity and is now standard practice for qualifying powder-bed parts; HIPed LPBF CoCrFeMnNi retains essentially the same chloride corrosion performance as the as-built state while improving density \cite{lpbf_cantor_2024}.

\subsection{Post-processing: the heat-treatment trade-off}

Heat treatment of AM HEAs is a genuine design decision with mechanical benefits and corrosion costs. Homogenization anneals relieve residual stress and coarsen the microstructure, but they can also restore the elemental segregation that rapid solidification suppressed. The clearest evidence comes from SLM AlCoCrFeNi$_{2.1}$: the as-built alloy, with its homogeneous inter-phase chemistry, shows superior corrosion resistance in 3.5~wt\% NaCl, while annealing dramatically degrades passivity as FCC/B2 segregation develops \cite{slm_ehea_2025}. Similar homogenization--corrosion interactions are reported for Al$_x$CoCrFeNi systems \cite{han2025}. Conversely, targeted annealing of LDED AlMo$_{0.25}$FeCoCrNi$_{2.1}$ has been used to tune strength and corrosion jointly \cite{lded_almofeco2024}. The design rule: any post-build heat treatment must be qualified for corrosion, not just mechanics.

\subsection{Representative mechanical performance}

A few results anchor the current state of the art. LPBF CoCrFeMnNi achieves strength above the wrought baseline with retained ductility, and exhibits notable abrasion resistance---the first low-stress abrasion data reported for any AM HEA \cite{lpbf_cantor_2024}. Mo-alloyed EHEAs reach yield strengths near 926~MPa with compressive strains approaching 40\% \cite{mo_ehea_2026}. WAAM MoNbTaWTi deposits remain hard (around 533~HV) and maintain compressive strength from 500 to 1000~$^{\circ}$C, positioning them for hot, abrasive niches such as exhaust-side marine hardware \cite{liu2021waam}. WAAM Cantor-type walls deliver competitive tensile behavior with anisotropy tied to the columnar grain structure \cite{andersson2026waam}. Finally, hydrogen permeation through LPBF CoCrFeMnNi is markedly lower than through 316L despite a higher diffusion coefficient, a favorable sign for hydrogen uptake and embrittlement resistance in cathodically protected seawater systems \cite{lpbf_cantor_2024}.
